%% ============================================================
%%  SECCIÓN XI — Análisis de Gráficas
%% ============================================================

\section{Análisis de Gráficas}
\label{sec:analisis_graficas}

\subsection{Gráfica 1: \texorpdfstring{$F$ vs.\ $\Delta\ell$}{F vs. Delta l} — Resorte Metálico}

La Fig.~\ref{fig:F_resorte} muestra una relación esencialmente lineal entre
la fuerza aplicada $F$ y la deformación longitudinal $\Delta\ell$ del
resorte metálico. El coeficiente de determinación $R^2 = 0{,}9966$ confirma
que el modelo lineal describe con alta fidelidad el comportamiento del
material dentro del rango experimental ensayado, lo que es consistente con la
Ley de Hooke. La pendiente del ajuste entrega directamente la constante
elástica del resorte:

\begin{equation*}
    k = \qty{80{,}89 \pm 2{,}38}{N/m}
\end{equation*}

Este valor caracteriza la rigidez del resorte y concuerda con el obtenido
mediante el ajuste por mínimos cuadrados de la
Sección~\ref{sec:tratamiento}. La ausencia de desviaciones sistemáticas
respecto a la recta de ajuste indica que el resorte permaneció dentro de
su región elástica, sin evidencia de deformación plástica ni de saturación.

\subsection{Gráfica 2: \texorpdfstring{$\sigma$ vs.\ $\varepsilon$}{sigma vs. epsilon} — Resorte Metálico}

La Fig.~\ref{fig:sigma_resorte} presenta el diagrama de esfuerzo real--deformación
del resorte. Al igual que en la Gráfica~1, el ajuste lineal resulta
altamente satisfactorio ($R^2 = 0{,}9971$), confirmando que el resorte
exhibe un comportamiento elástico lineal. La pendiente de la recta corresponde 
al módulo de Young efectivo recalculado mediante el área instantánea:

\begin{equation*}
    Y_r = \qty{90{,}23 \pm 2{,}41}{kPa}
\end{equation*}

Este valor, notablemente inferior al del acero masivo
($Y_\text{acero} \sim \qty{200}{GPa}$), no representa una propiedad
intrínseca del material sino un módulo efectivo macroscópico del sistema
resorte. La discrepancia se explica por la geometría helicoidal: bajo carga 
axial, el alambre experimenta predominantemente torsión y flexión, lo que 
reduce el módulo efectivo en varios órdenes de magnitud respecto a una 
barra sólida.

\subsection{Gráfica 3: \texorpdfstring{$F$ vs.\ $\Delta\ell$}{F vs. Delta l} — Liga Elástica (Carga y Descarga)}

La Fig.~\ref{fig:F_liga} permite visualizar el comportamiento crudo de la liga 
elástica. Se observa que la curva de descarga no coincide con la de carga, 
presentando una elongación mayor para un mismo valor de fuerza durante el 
retorno. Esta separación evidencia cualitativamente la existencia de histéresis 
y una recuperación lenta de la estructura interna, independientemente de las 
variaciones en la sección transversal.

\subsection{Gráfica 4: \texorpdfstring{$\sigma$ vs.\ $\varepsilon$}{sigma vs. epsilon} — Liga Elástica (Carga)}

La Fig.~\ref{fig:sigma_liga_carga} muestra el diagrama de esfuerzo real para 
la liga. A diferencia del resorte, la pendiente aumenta con la deformación, 
un endurecimiento característico atribuido a la alineación de cadenas 
poliméricas. Al emplear el esfuerzo verdadero ($\sigma = F/S$), la curvatura 
es más pronunciada debido a que el área transversal disminuye significativamente 
durante el estiramiento.

Aunque se calculó un módulo de Young de referencia ($Y_l \approx \qty{538}{kPa}$), 
se destaca que este valor representa únicamente un \textbf{módulo secante promedio} 
sobre la región ensayada y carece de un sentido físico global dado el carácter 
manifiestamente no lineal del material. El modelo polinomial de grado~2 resulta 
mucho más representativo:

\begin{equation*}
    \sigma(\varepsilon) = 340{,}20\,\varepsilon^2 + 372{,}28\,\varepsilon
                        + 38{,}82 \quad \text{(kPa)}
\end{equation*}

\subsection{Gráfica 5: \texorpdfstring{$\sigma$ vs.\ $\varepsilon$}{sigma vs. epsilon} — Liga Elástica (Descarga)}

La Fig.~\ref{fig:sigma_liga_descarga} muestra la trayectoria de recuperación. 
El ajuste polinomial de grado~2 modela adecuadamente el retorno ($R^2 = 0{,}9667$):

\begin{equation*}
    \sigma(\varepsilon) = 3119{,}97\,\varepsilon^2 - 1024{,}06\,\varepsilon
                        + 154{,}01 \quad \text{(kPa)}
\end{equation*}

La deformación residual observada durante la descarga confirma que la liga 
no restituye de forma inmediata la energía elástica. Este proceso es 
termodinámicamente irreversible y las fuerzas involucradas son no conservativas, 
lo cual se cuantifica en la Tabla~\ref{tab:area_histeresis}.

\subsection{Gráfica 6: Lazo de histéresis — Liga Elástica}

La Fig.~\ref{fig:histeresis_lazo} muestra el lazo de histéresis cerrado. 
El área encerrada por el lazo, calculada mediante la integración numérica de la 
diferencia entre las curvas de carga y descarga en la 
Sección~\ref{sec:tratamiento}, es:

\begin{equation*}
    A = \qty{13{,}817}{kJ/m^3}
\end{equation*}

Este valor representa la energía mecánica disipada internamente por unidad de 
volumen en forma de calor. La asimetría del lazo es consistente con la 
viscoelasticidad del material, donde la fricción intramolecular alcanza su 
máximo en el rango de deformaciones intermedias.